\subsection*{CU-25: Reportar a un jugador}
\addcontentsline{toc}{subsection}{CU-25: Reportar a un jugador}
\label{cu-25}

\begin{fichacu}{CU-25}{Reportar a un jugador}
  \crudFila{Responsable}{José Eduardo Prior Hernández}
  \crudFila{Actor principal}{Jugador}
  \crudFila{Actores de sistema}{Servidor de partidas, Base de datos}
  \crudFila{Prototipos}{10a, 10b}
  \crudFila{Prioridad}{Must (debe estar en la entrega). Categoría (a) Obligatorio: característica 7 del cliente (jugadores expulsados) y moderación por reportes (\mbox{CON-09}): los strikes restringen, suspenden o expulsan al jugador sin intervención de moderadores.}
  \crudFila{Objetivo}{Denunciar la conducta de otro jugador, adjuntando automáticamente las pruebas de la partida o de la sala, y sancionar de forma automática al jugador que reúne reportes de varios jugadores distintos, con un strike cuya sanción crece con cada uno.}
  \crudFila{Disparador}{El Jugador selecciona ``Reportar'' en el menú de un jugador, accesible desde el chat, la partida, la sala de espera o el perfil del jugador (\mbox{CU-06}).}
  \textbf{Precondiciones} & \cuCelda{\begin{cureglas}
      \item \textbf{\mbox{PRE-01}} El Jugador tiene una \textit{Sesion} abierta y su token es válido.
      \item \textbf{\mbox{PRE-02}} El \textit{Usuario} tiene \texttt{tipo\_\allowbreak{}cuenta} en \texttt{REGISTRADA} (\mbox{RN-07}). \textit{(Ver \mbox{FA-01})}
      \item \textbf{\mbox{PRE-03}} El jugador reportado existe y no es el propio denunciante.
      \item \textbf{\mbox{PRE-04}} Los diccionarios de recursos contienen el texto de los cinco motivos de reporte.
      \item \textbf{\mbox{PRE-05}} El Servidor de partidas está en ejecución y tiene conexión disponible con la Base de datos.
    \end{cureglas}} \\ \hline
  \textbf{Flujo normal} & \cuCelda{\begin{cuenum}[start=1]
      \item El Jugador abre el menú de un jugador y selecciona ``Reportar''.
      \item El SJM despliega la \texttt{GUI\_\allowbreak{}Report} con los cinco motivos, un campo de descripción opcional y el aviso de qué pruebas se adjuntarán (\mbox{RN-02}).
      \item El Jugador elige un motivo, escribe la descripción si quiere y da click en ``Enviar reporte''. \textit{(Ver \mbox{FA-02}, \mbox{FA-03})}
      \item El SJM valida que haya un motivo elegido y que la descripción no exceda el largo máximo. \textit{(Ver \mbox{FA-04})}
      \item El SJM envía el mensaje \texttt{REPORT\_\allowbreak{}PLAYER\_\allowbreak{}REQUEST} con el reportado, el motivo, la descripción, la partida o la sala desde la que se reporta si la hay, y el token de sesión. \textit{(Ver \mbox{EX-02}, \mbox{EX-03})}
      \item El Servidor de partidas verifica que el token corresponda a una \textit{Sesion} abierta de una cuenta registrada (\mbox{RN-07}). \textit{(Ver \mbox{FA-01}, \mbox{FA-05})}
      \item El Servidor de partidas verifica que el reportado exista, no sea el denunciante y no tenga \texttt{estado\_\allowbreak{}cuenta} en \texttt{ELIMINADA}; si se adjunta una partida, que denunciante y reportado tengan \textit{Participacion} en ella, y si se adjunta una sala, que ambos ocupen una plaza en ella (\mbox{RN-09}). \textit{(Ver \mbox{FA-06})}
      \item El Servidor de partidas verifica que el \texttt{motivo} sea uno de los cinco admitidos. \textit{(Ver \mbox{FA-07})}
    \end{cuenum}} \\
   & \cuCelda{\begin{cuenum}[start=9]
      \item El Servidor de partidas verifica que el denunciante no haya reportado ya al mismo jugador por la misma partida o la misma sala (\mbox{RN-03}). \textit{(Ver \mbox{FA-08})}
      \item El Servidor de partidas comprueba cuántos reportes ha emitido el denunciante en el periodo vigente (\mbox{RN-06}). \textit{(Ver \mbox{FA-09})}
      \item El Servidor de partidas inicia una transacción sobre la Base de datos, aislada de modo que dos reportes simultáneos contra el mismo jugador no cuenten dos veces los mismos reportes para un strike.
      \item El Servidor de partidas crea el \textit{Reporte} con el denunciante, el reportado, el motivo, la descripción, la fecha y el estado \texttt{PENDIENTE}. \textit{(Ver \mbox{EX-01})}
      \item Si el reporte procede de una partida, el Servidor de partidas asocia al \textit{Reporte} la \textit{Partida}, y si procede de una sala, la \textit{Sala}, de modo que las jugadas y el chat queden adjuntos por referencia y no por copia (\mbox{RN-01}).
      \item Si la cuenta del reportado no está \texttt{BANEADA}, el Servidor de partidas cuenta los denunciantes distintos de sus \textit{Reporte} que siguen en \texttt{PENDIENTE}, incluido el que acaba de crear (\mbox{RN-10}, \mbox{RN-12}). \textit{(Ver \mbox{FA-10})}
      \item El Servidor de partidas confirma la transacción. \textit{(Ver \mbox{EX-04})}
      \item El Servidor de partidas responde con \texttt{REPORT\_\allowbreak{}OK}.
    \end{cuenum}} \\
   & \cuCelda{\begin{cuenum}[start=17]
      \item El SJM muestra la \texttt{GUI\_\allowbreak{}MessageSuccess} indicando que el reporte se envió, sin informar de su resultado ni de si el reportado fue sancionado (\mbox{RN-05}).
      \item Fin del caso de uso.
    \end{cuenum}} \\ \hline
  \textbf{Flujos alternos} & \cuCelda{\textbf{\mbox{FA-01}: El Jugador es un invitado}\par
    \begin{cuenum}[start=1]
      \item El SJM no ofrece la opción de reportar e indica que reportar requiere una cuenta registrada (\mbox{RN-07}).
      \item Si aun así llega una solicitud de una cuenta de invitado, el Servidor de partidas la rechaza en el paso 6 con \texttt{REPORT\_\allowbreak{}FAILED} y el motivo \texttt{CUENTA\_\allowbreak{}INVITADO}, y registra el intento en la bitácora del servidor como indicio de cliente modificado.
      \item Fin del caso de uso.
    \end{cuenum}} \\
   & \cuCelda{\textbf{\mbox{FA-02}: El Jugador cancela el reporte}\par
    \begin{cuenum}[start=1]
      \item El Jugador selecciona ``Cancelar''.
      \item El SJM cierra la \texttt{GUI\_\allowbreak{}Report} sin enviar nada.
      \item Fin del caso de uso.
    \end{cuenum}} \\
   & \cuCelda{\textbf{\mbox{FA-03}: El Jugador reporta desde fuera de una partida o una sala}\par
    \begin{cuenum}[start=1]
      \item El SJM no adjunta ninguna \textit{Partida} ni ninguna \textit{Sala} al reporte, porque no hay ninguna de la que hablar; por ejemplo, si reporta desde el perfil del \mbox{CU-06}.
      \item El SJM indica en el formulario que no se adjuntarán pruebas y que la descripción es, por tanto, lo único que acompañará al reporte.
      \item El flujo continúa en el paso 4 del FN.
    \end{cuenum}} \\
   & \cuCelda{\textbf{\mbox{FA-04}: Falta el motivo o la descripción excede el largo}\par
    \begin{cuenum}[start=1]
      \item El SJM resalta el campo correspondiente y escribe junto a él qué se espera.
      \item El flujo continúa en el paso 3 del FN.
    \end{cuenum}} \\
   & \cuCelda{\textbf{\mbox{FA-05}: La sesión ya no es válida}\par
    \begin{cuenum}[start=1]
      \item El Servidor de partidas responde con \texttt{SESSION\_\allowbreak{}INVALID}.
      \item El SJM muestra la \texttt{GUI\_\allowbreak{}MessageWarning} y despliega la \texttt{GUI\_\allowbreak{}Login}.
      \item Fin del caso de uso.
    \end{cuenum}} \\
   & \cuCelda{\textbf{\mbox{FA-06}: El reportado no existe o es el propio denunciante}\par
    \begin{cuenum}[start=1]
      \item El Servidor de partidas responde con \texttt{REPORT\_\allowbreak{}FAILED} y el motivo \texttt{DESTINATARIO\_\allowbreak{}INVALIDO}.
      \item El Servidor de partidas registra el intento en la bitácora del servidor si el reportado es el propio denunciante, por tratarse de un indicio de cliente modificado.
      \item El SJM muestra la \texttt{GUI\_\allowbreak{}MessageWarning}.
      \item Fin del caso de uso.
    \end{cuenum}} \\
   & \cuCelda{\textbf{\mbox{FA-07}: El motivo no pertenece al catálogo}\par
    \begin{cuenum}[start=1]
      \item El Servidor de partidas responde con \texttt{REPORT\_\allowbreak{}FAILED} y el motivo \texttt{MOTIVO\_\allowbreak{}INVALIDO}, y registra la discrepancia en la bitácora del servidor.
      \item El SJM recarga el catálogo de motivos.
      \item El flujo continúa en el paso 2 del FN.
    \end{cuenum}} \\
   & \cuCelda{\textbf{\mbox{FA-08}: Ya existe un reporte del mismo denunciante sobre la misma partida o sala}\par
    \begin{cuenum}[start=1]
      \item El Servidor de partidas responde con \texttt{REPORT\_\allowbreak{}FAILED} y el motivo \texttt{REPORTE\_\allowbreak{}DUPLICADO} (\mbox{RN-03}).
      \item El SJM indica que ya reportó a ese jugador por esa partida o esa sala.
      \item Fin del caso de uso.
    \end{cuenum}} \\
   & \cuCelda{\textbf{\mbox{FA-09}: El Jugador excedió el número de reportes permitido}\par
    \begin{cuenum}[start=1]
      \item El Servidor de partidas responde con \texttt{REPORT\_\allowbreak{}THROTTLED} y el tiempo que debe esperar.
      \item El SJM lo indica mediante la \texttt{GUI\_\allowbreak{}MessageWarning}.
      \item Fin del caso de uso.
    \end{cuenum}} \\
   & \cuCelda{\textbf{\mbox{FA-10}: El reportado reúne reportes de tres jugadores distintos}\par
    \begin{cuenum}[start=1]
      \item El Servidor de partidas encuentra al menos tres denunciantes distintos entre los \textit{Reporte} pendientes del reportado (\mbox{RN-10}) y, por las \textit{Sancion} que ya recibió, determina qué número de strike le corresponde (\mbox{RN-12}).
      \item El Servidor de partidas crea la \textit{Sancion} de ese strike, aplicada por el sistema y sin moderador (\mbox{RN-04}, \mbox{RN-11}): de ámbito \texttt{CHAT} y veinticuatro horas para el primero, de ámbito \texttt{CUENTA} y siete días para el segundo, y de ámbito \texttt{CUENTA} sin \texttt{fecha\_\allowbreak{}fin} para el tercero. \textit{(Ver \mbox{EX-01})}
      \item Si la \textit{Sancion} es de ámbito \texttt{CUENTA}, el Servidor de partidas cambia el \texttt{estado\_\allowbreak{}cuenta} del reportado a \texttt{SUSPENDIDA} o \texttt{BANEADA}, según el strike, y cierra todas sus \textit{Sesion} abiertas con el \texttt{motivo\_\allowbreak{}cierre} \texttt{SANCION} (\mbox{RN-13}).
      \item El Servidor de partidas pasa a \texttt{ATENDIDO} todos los \textit{Reporte} pendientes del reportado y los asocia a la \textit{Sancion} (\mbox{RN-10}).
      \item El Servidor de partidas registra la sanción en la \textit{BitacoraModeracion} con la acción \texttt{SANCION\_\allowbreak{}APLICADA}, el reportado como usuario afectado y, en el detalle, el número de strike y los reportes que lo originaron.
      \item El Servidor de partidas confirma la transacción. \textit{(Ver \mbox{EX-04})}
      \item Si la \textit{Sancion} es de ámbito \texttt{CUENTA}, el Servidor de partidas envía \texttt{ACCOUNT\_\allowbreak{}SANCTIONED} a los SJM conectados del sancionado, que despliegan la \texttt{GUI\_\allowbreak{}BannedAccount}, y lo saca de la sala o de la partida en la que esté (\mbox{RN-13}). \textit{(Ver \mbox{EX-05})}
      \item Si la \textit{Sancion} es de ámbito \texttt{CHAT}, el Servidor de partidas envía \texttt{CHAT\_\allowbreak{}RESTRICTED} a los SJM conectados del sancionado, que deshabilitan la escritura del chat conforme al \mbox{FA-03} del \mbox{CU-18}.
      \item El flujo continúa en el paso 16 del FN.
    \end{cuenum}} \\ \hline
  \textbf{Excepciones} & \cuCelda{\textbf{\mbox{EX-01}: Fallo al escribir en la base de datos}\par
    \begin{cuenum}[start=1]
      \item El Servidor de partidas revierte la transacción, de modo que no quede un \textit{Reporte} sin su partida o su sala asociada ni un strike aplicado a medias.
      \item El Servidor de partidas registra la excepción con un identificador de incidencia y responde con \texttt{SERVER\_\allowbreak{}ERROR}.
      \item El SJM muestra la \texttt{GUI\_\allowbreak{}MessageError} con el identificador y conserva lo capturado para que pueda reintentarlo.
      \item El flujo continúa en el paso 3 del FN.
    \end{cuenum}} \\
   & \cuCelda{\textbf{\mbox{EX-02}: Se agota el tiempo de espera de la respuesta}\par
    \begin{cuenum}[start=1]
      \item El SJM no reenvía el reporte por su cuenta: podría duplicarlo.
      \item El SJM muestra la \texttt{GUI\_\allowbreak{}MessageError} indicando que el reporte pudo haberse enviado y que, si ya se guardó, un nuevo intento se rechazará como duplicado (\mbox{FA-08}).
      \item Fin del caso de uso.
    \end{cuenum}} \\
   & \cuCelda{\textbf{\mbox{EX-03}: La conexión se interrumpe}\par
    \begin{cuenum}[start=1]
      \item El SJM muestra la barra de conexión perdida y conserva lo capturado.
      \item Al recuperar la conexión, el flujo continúa en el paso 3 del FN.
    \end{cuenum}} \\
   & \cuCelda{\textbf{\mbox{EX-04}: Fallo al confirmar la transacción}\par
    \begin{cuenum}[start=1]
      \item El Servidor de partidas trata la transacción como fallida y registra la excepción señalando que el estado del reporte, y del strike si lo había, es indeterminado.
      \item El SJM muestra la \texttt{GUI\_\allowbreak{}MessageError} indicando que el reporte pudo haberse enviado y que, si ya se guardó, un nuevo intento se rechazará como duplicado (\mbox{FA-08}).
      \item Fin del caso de uso.
    \end{cuenum}} \\
   & \cuCelda{\textbf{\mbox{EX-05}: No puede retirarse al sancionado de su partida}\par
    \begin{cuenum}[start=1]
      \item El Servidor de partidas registra la excepción con un identificador de incidencia. La \textit{Sancion} ya está confirmada y no se revierte.
      \item Como el sancionado ya no tiene sesiones abiertas, la partida lo trata como desconectado (\mbox{CU-17}) y su rival puede reclamar la victoria a los sesenta segundos (\mbox{CU-15}).
      \item El flujo continúa en el paso 16 del FN.
    \end{cuenum}} \\ \hline
  \textbf{Postcondiciones} & \cuCelda{\begin{cureglas}
      \item \textbf{\mbox{POST-01}} Si el caso termina por el FN sin pasar por el \mbox{FA-10}, existe un \textit{Reporte} con estado \texttt{PENDIENTE}, su denunciante, su reportado, su motivo y su fecha.
      \item \textbf{\mbox{POST-02}} Si el caso termina por el FN y el reporte procede de una partida, el \textit{Reporte} referencia esa \textit{Partida}, y por tanto sus \textit{Jugada} y sus \textit{Mensaje}; si procede de una sala, referencia esa \textit{Sala} y sus \textit{Mensaje}.
      \item \textbf{\mbox{POST-03}} Si el caso termina por el FN sin pasar por el \mbox{FA-10}, la cuenta del reportado \textbf{no} cambió: un reporte por sí solo no sanciona (\mbox{RN-04}).
      \item \textbf{\mbox{POST-04}} Si el caso termina por el FN, el denunciante no puede volver a reportar al mismo jugador por la misma partida o la misma sala.
      \item \textbf{\mbox{POST-05}} No puede existir más de un \textit{Reporte} del mismo denunciante sobre el mismo reportado y la misma partida o sala.
      \item \textbf{\mbox{POST-06}} Si el caso termina por cualquier flujo alterno o excepción, no se creó ningún \textit{Reporte}.
      \item \textbf{\mbox{POST-07}} Si el caso pasa por el \mbox{FA-10}, el reportado tiene una \textit{Sancion} nueva sin moderador, con el ámbito y la duración de su número de strike; ninguno de sus \textit{Reporte} sigue en \texttt{PENDIENTE}, y la sanción consta en la \textit{BitacoraModeracion}.
      \item \textbf{\mbox{POST-08}} Si la \textit{Sancion} del \mbox{FA-10} es de ámbito \texttt{CUENTA}, el reportado tiene \texttt{estado\_\allowbreak{}cuenta} en \texttt{SUSPENDIDA} o \texttt{BANEADA}, no tiene ninguna \textit{Sesion} abierta y no ocupa plaza en ninguna sala ni sigue en ninguna partida en curso, salvo en el \mbox{EX-05}.
    \end{cureglas}} \\ \hline
  \textbf{Reglas de negocio} & \cuCelda{\begin{cureglas}
      \item \textbf{\mbox{RN-01}} El reporte adjunta las jugadas y el chat de la partida, o el chat de la sala, de forma automática, por referencia a la \textit{Partida} o a la \textit{Sala}. No se copian los mensajes ni las jugadas: se apunta a ellos, de modo que el reporte no duplique información que ya existe y no pueda contradecirla.
      \item \textbf{\mbox{RN-02}} El Jugador ve, antes de enviar, qué pruebas se adjuntarán. Adjuntar su chat sin decírselo sería recoger datos a su espalda.
      \item \textbf{\mbox{RN-03}} Un denunciante no puede reportar dos veces al mismo jugador por la misma partida ni por la misma sala. Sí puede hacerlo por partidas o salas distintas. Un reporte sin partida ni sala se trata como duplicado si ya existe otro sin partida ni sala del mismo denunciante al mismo jugador en las últimas veinticuatro horas.
      \item \textbf{\mbox{RN-04}} Un reporte por sí solo no sanciona. La única consecuencia automática de reportar es el strike del \mbox{RN-10}, que aplica el sistema sin la revisión de un moderador y sin apelación.
      \item \textbf{\mbox{RN-05}} Al denunciante no se le informa del resultado del reporte ni de si el reportado fue sancionado.
      \item \textbf{\mbox{RN-06}} El número de reportes que un jugador puede emitir en un periodo está acotado, para que reportar no sirva como forma de hostigamiento.
      \item \textbf{\mbox{RN-07}} Las cuentas de invitado no pueden reportar, pero sí ser reportadas. Un reporte sin cuenta rastreable no puede sostener el \mbox{RN-06} ni responder de un uso abusivo, y como una cuenta de invitado se crea sin correo, una sola persona podría reunir con ellas los tres denunciantes de un strike. El invitado, en cambio, juega y escribe en el chat como cualquier otro, así que su conducta sí puede reportarse.
    \end{cureglas}} \\
   & \cuCelda{\begin{cureglas}
      \item \textbf{\mbox{RN-08}} Hay cinco motivos de reporte: \texttt{LENGUAJE\_\allowbreak{}OFENSIVO}, \texttt{ACOSO}, \texttt{NOMBRE\_\allowbreak{}INAPROPIADO}, \texttt{TRAMPAS} y \texttt{JUEGO\_\allowbreak{}ANTIDEPORTIVO}. Son un dominio cerrado de \textit{Reporte}, declarado con \texttt{CHECK}, y no se cambian desde la aplicación (\mbox{D-23}).
      \item \textbf{\mbox{RN-09}} Solo puede adjuntarse una partida en la que participaron el denunciante y el reportado, o la sala en la que ambos ocupan una plaza al reportar. De otro modo el reporte adjuntaría pruebas ajenas a los hechos.
      \item \textbf{\mbox{RN-10}} Cuando un jugador reúne \textit{Reporte} en \texttt{PENDIENTE} de tres denunciantes distintos, el sistema le aplica un strike en la misma transacción en la que se crea el último y pasa todos sus reportes pendientes a \texttt{ATENDIDO}, de modo que ninguno vuelva a contar para otro strike.
      \item \textbf{\mbox{RN-11}} Cada strike tiene una sanción fija. El primero restringe el chat veinticuatro horas (\textit{Sancion} de ámbito \texttt{CHAT}). El segundo suspende la cuenta siete días (\textit{Sancion} de ámbito \texttt{CUENTA}, \texttt{estado\_\allowbreak{}cuenta} en \texttt{SUSPENDIDA}). El tercero la expulsa de forma permanente (\textit{Sancion} de ámbito \texttt{CUENTA} sin \texttt{fecha\_\allowbreak{}fin}, \texttt{estado\_\allowbreak{}cuenta} en \texttt{BANEADA}).
      \item \textbf{\mbox{RN-12}} Los strikes no caducan en esta versión. El número de un strike es el de sanciones que el jugador ya recibió, más uno. Una cuenta \texttt{BANEADA} no recibe más strikes: los reportes posteriores contra ella se guardan sin consecuencia.
      \item \textbf{\mbox{RN-13}} Una sanción de ámbito \texttt{CUENTA} surte efecto de inmediato: cierra todas las sesiones del sancionado (\mbox{RN-07} del \mbox{CU-01}), lo saca de la sala en la que espere como si hubiera salido de ella y, si está jugando, lo retira de la partida con las consecuencias de un abandono (\mbox{CU-15}). Una de ámbito \texttt{CHAT} solo le impide escribir en el chat (\mbox{RN-05} del \mbox{CU-18}).
    \end{cureglas}} \\ \hline
  \textbf{Observación} & \cuCelda{\textit{Sobre por qué las pruebas van por referencia.}\par
    Se consideró copiar en el reporte los mensajes y las jugadas relevantes, para que la prueba quedara congelada en el momento de reportar. Se descartó por dos motivos. El primero es que duplicar el chat en el reporte crea dos copias del mismo texto que pueden divergir. El segundo es que recortar qué mensajes son relevantes obliga a decidirlo en el momento del envío, cuando todavía no se sabe qué hará falta mirar. Referenciar la partida o la sala entrega el contexto completo, que es lo que después explica por qué se aplicó cada strike.} \\ \hline
  \textbf{Observación} & \cuCelda{\textit{Sobre por qué los strikes son automáticos.}\par
    Se consideró que un moderador revisara cada reporte antes de sancionar, porque distingue mejor un reporte fundado de uno malintencionado. Se descartó porque el cliente pidió una moderación simple (\mbox{CON-09}): la revisión humana obliga a construir el rol de moderador, su panel y su cola de reportes, y deja la conducta sin consecuencia mientras nadie la atiende. Exigir tres denunciantes distintos sustituye ese juicio por un acuerdo mínimo entre jugadores, y la escala de strikes hace que el primero cueste poco: solo restringe el chat un día.} \\ \hline
  \textbf{Datos que toca} & \cuCelda{\begin{cudatos}
      \item \textbf{\textit{Sesion}:} lee por token \textit{(paso 6)}; cierra las abiertas del sancionado \textit{(\mbox{FA-10})}
      \item \textbf{\textit{Usuario}:} lee el \texttt{tipo\_\allowbreak{}cuenta} del denunciante y el estado del reportado \textit{(pasos 6, 7, 14)}; escribe el \texttt{estado\_\allowbreak{}cuenta} del sancionado \textit{(\mbox{FA-10})}
      \item \textbf{\textit{Participacion}, \textit{SalaParticipante}:} lee, para comprobar que ambos jugaron la partida adjunta o están en la sala adjunta \textit{(paso 7)}
      \item \textbf{Motivos de reporte:} el dominio de \texttt{Reporte.motivo} \textit{(pasos 2, 8, \mbox{D-23})}
      \item \textbf{\textit{Reporte}:} lee, para comprobar duplicados, el límite y los pendientes del reportado \textit{(pasos 9, 10, 14)}
      \item \textbf{\textit{Reporte}:} crea \textit{(pasos 12, 13)}; pasa los pendientes del reportado a \texttt{ATENDIDO} \textit{(\mbox{FA-10})}
      \item \textbf{\textit{Partida}, \textit{Sala}:} referencia, sin copiar nada de ellas \textit{(paso 13)}
      \item \textbf{\textit{Sancion}:} lee las del reportado y crea la del strike \textit{(\mbox{FA-10})}
      \item \textbf{\textit{BitacoraModeracion}:} crea \textit{(\mbox{FA-10})}
      \item \textbf{\textit{Participacion}, \textit{EstadisticaModo}:} las que escribe el \mbox{CU-15} al retirar al sancionado de su partida \textit{(\mbox{FA-10})}
    \end{cudatos}} \\ \hline
\end{fichacu}
\clearpage
